\section{Complete fleets and marginal-price support}\label{sec:model}

The unit of optimization is a \emph{complete} fleet plan, rather than an
individual vehicle duty.  Fix a timetable, allowed directed movements, a
charging horizon, and the equipment and energy assumptions.  Let
$\Fleet\ne\varnothing$ contain precisely the plans that cover every mandatory
service once, assign the services to used buses, route those buses through
allowed movements, respect battery and charger constraints, and restore every
used bus to its required terminal inventory.  For $x\in\Fleet$, write $c(x)$
for intrinsic operating cost and $e(x)\in\R_+^T$ for grid energy bought in the
$T$ market intervals.  An unused bus contributes neither charging nor cost.
We use the same $\Fleet$, $c$, and $e$ for physical planning, price response,
and convexification.  Finitely many discrete duties with bounded continuous
charging give a compact $\Fleet$ in the cases below; assume more generally
that its image $\{(e(x),c(x)):x\in\Fleet\}$ is compact.

In the timed physical model, each used path starts at battery inventory $B$.
A service or directed movement consumes its declared energy in chronological
order.  Inventory remains between reserve $r$ and capacity $B$ at every
event.  A depot charging session buying $q$ grid kWh adds $\eta q$ inventory,
where $0<\eta\leq1$, within its availability window and bus-power limit.
Concurrent sessions obey the shared power or connector limit.  Each path
finishes at inventory $B$ by the declared horizon.  In the homogeneous
single-connector formulation, charging intervals are split at all visit and
market boundaries.  If interval $k$ has duration $d_k$ and power $P_k$, its
charges obey $0\leq z_{mk}\leq P_kd_k y_m$ and
$\sum_m z_{mk}\leq P_kd_k$, where $y_m$ indicates the selected visit.
Serial allocation of $z_{mk}/P_k$ hours to the visits realizes these totals
when switching time is zero, efficiency is constant, and charging does not
taper.  The worked constructions below assume zero travel cost and no
unlisted travel or auxiliary load; the public case declares its movement
energy separately.

Let $F:\R^T\to\R\cup\{+\infty\}$ be a proper, closed, convex supply cost,
finite at every physical load.  Its effective domain includes the feasible
loads but may also encode generation capacity and ramp limits.  Write
$\mathcal E=\operatorname{conv}\{e(x):x\in\Fleet\}$.  For the support
equivalence, assume either
$\operatorname{ri}\mathcal E\cap\operatorname{ri}(\operatorname{dom}F)
\ne\varnothing$, or that $F$ has a finite convex minorant on $\R^T$ that
agrees with it on $\mathcal E$.  The latter covers the nonnegative-domain
quadratic used below via its unrestricted quadratic formula.
Here $F$ represents the supplier's aggregate incremental provision cost for
fleet charging, conditional on background demand already accounted for. A
system operator posts the marginal tariff calculated from a candidate or
forecast fleet load before the fleet chooses a schedule. The tariff is held
fixed when evaluating alternative complete fleets. This ex-ante pricing rule
is the institution tested below; the supplier is not a strategic player and
the fleet does not anticipate a tariff reset during its deviation.
The physical and complete-fleet-hull values are
\begin{equation}\label{eq:theory-values}
\begin{aligned}
 D&=\min_{x\in\Fleet}\{c(x)+F(e(x))\},\\
 CH&=\min_{(e,c)\in\operatorname{conv}\{(e(x),c(x)):x\in\Fleet\}}
       \{c+F(e)\},\qquad \Delta=D-CH\geq0.
\end{aligned}
\end{equation}
The hull mixes whole feasible plans, including their charger constraints;
its mean load is a pricing and bounding object, not necessarily a dispatch.
For a posted price $p$, define the complete-fleet response
$V(p)=\min_{x\in\Fleet}\{c(x)+p^\top e(x)\}$.  A marginal price at a
physical plan $x$ is any $p\in\partial F(e(x))$; this set can contain more
than one price and need not be nonempty at every boundary load.  Define
\begin{equation}\label{eq:theory-regret}
 r(x;p)=c(x)+p^\top e(x)-V(p).
\end{equation}
For a stated price selection $p_x$, write $r(x)=r(x;p_x)$; a nonsmooth
model must declare that selection when reporting own-price regret.
The alternative in $V(p)$ takes the posted price as fixed, as the stated
tariff rule requires.  A strategic fleet instead compares bills such as
$c(y)+p_y^\top e(y)$ across its alternatives $y$ and anticipates
the resulting price change.  That objective differs from both
$c(y)+p^\top e(y)$ and the planner's $c(y)+F(e(y))$; it needs a specified
settlement rule and strategic game.  The present price-support and
compensation statements make no claim about strategic behavior
\citep{oneill2005,gribikHoganPope2007,ganTopcuLow2013}.

\begin{proposition}[Complete-fleet support]\label{prop:theory-support}
Under the common-model and attainment assumptions above, every physical
$x\in\Fleet$ and every $p\in\partial F(e(x))$ satisfy
\begin{equation}\label{eq:theory-support}
 r(x;p)\geq c(x)+F(e(x))-CH\geq D-CH=\Delta.
\end{equation}
Moreover, $D=CH$ if and only if some physical plan is a best response at
some price in its subdifferential.  If $D=CH$, every physical planner
optimizer is supported by every attained maximizer of the hull dual
$V(p)-F^*(p)$.
\end{proposition}

The proof in Appendix~\ref{app:theory-proofs} is a specialization of classical
nonconvex price-support and convexification arguments
\citep{starr1969,oneill2005,gribikHoganPope2007,huaBaldick2017,schiro2016}.
It says nothing about whether an iterative price adjustment converges or
which optimum an algorithm selects.  The stated qualification
ensures strong Fenchel duality and an attained supporting price.  If an
economic price convention admits only a subset of $\partial F(e(x))$, the
converse additionally requires an attained supporting price in that subset.
In particular, a normal caused solely by extending a nonnegative-load cost
as $+\infty$ to negative loads is not automatically a generation price.

The accounting identity also requires a stated deviation right.  In this
paper one fleet controls all buses and its alternative must be a member of
the same $\Fleet$.  Separate operators using shared chargers would need
reserved capacities, trades, or congestion rights before a unilateral
opportunity set could be defined.  For the single fleet, put
$F^*(p)=\sup_{u\in\R^T}\{p^\top u-F(u)\}$ and define
\begin{align}\label{eq:theory-loc}
 \mathrm{LOC}_{\rm fleet}(x;p)&=c(x)+p^\top e(x)-V(p),\\
 \mathrm{LOC}_{\rm supply}(x;p)&=F(e(x))-p^\top e(x)+F^*(p).
\end{align}
Whenever $F^*(p)$ is finite, both terms are nonnegative and their sum is
$c(x)+F(e(x))-[V(p)-F^*(p)]$.  At a hull-optimal price and a physical planner
optimizer, strong Fenchel duality makes that sum $\Delta$; the fleet and
supply portions depend on the chosen common price.  At any
$p\in\partial F(e(x))$, supply LOC is zero and fleet LOC equals $r(x;p)$.
No budget-balance or payment institution
follows from this arithmetic.

For quadratic supply $F(e)=a^\top e+\tfrac12 e^\top Qe$ on $\R_+^T$ with
$Q\succeq0$, take the unrestricted-quadratic gradient price
$p_x=a+Qe(x)$ and
write $\|z\|_Q=(z^\top Qz)^{1/2}$.  There is a useful complement
to the lower bound.  Let
$R_Q=\max_{x,y\in\Fleet}\|e(x)-e(y)\|_Q$ and
$H(x)=c(x)+F(e(x))-CH$.  Then
\begin{equation}\label{eq:theory-upper}
 H(x)\leq r(x)\leq H(x)+R_Q\sqrt{2H(x)}.
\end{equation}
Thus a small cost gap controls own-price regret only together with the
curvature-weighted range of feasible fleet loads.  The proof permits singular
$Q$ and boundary loads and appears in Appendix~\ref{app:theory-proofs}.

For numerical certification, valid bounds
$D\in[L_D,U_D]$ and $CH\in[L_{CH},U_{CH}]$ imply
$\Delta\in[L_D-U_{CH},U_D-L_{CH}]$.  A replay-feasible physical plan supplies
an upper bound on $D$; a mixture of replay-feasible complete plans supplies
an upper bound on $CH$.  At any price with finite $F^*(p)$, a global
price-response lower bound $\underline V(p)\leq V(p)$ gives the hull lower
bound $\underline V(p)-F^*(p)$.  A nearly solved restricted column pool alone
cannot supply that global lower bound.  In the separable implementation,
$F(e)=\sum_t(a_te_t+b_te_t^2/2)$ for $e\geq0$ and $b_t\geq0$;
the conjugate term is
$[p_t-a_t]_+^2/(2b_t)$ when $b_t>0$, and is finite at $b_t=0$ only for
$p_t\leq a_t$.  These are standard convexification and Fenchel certificates
over complete plans, rather than a new general convex-hull pricing method
\citep{huaBaldick2017}.  Section~\ref{sec:computational} compares two specified market states.
